% !TeX root = ../My_thesis.tex
\chapter{ПРОЕКТИРОВАНИЕ И РЕАЛИЗАЦИЯ МОДУЛЬНЫХ ТЕСТОВ} \label{ch3}

В данной главе описано проектирование и реализация модульных тестов для разработанного приложения фильтрации IP-адресов. В соответствии с заданием выбраны и применены четыре техники проектирования тестов: две на основе спецификации (Boundary Value Analysis и Equivalence Partitioning) и две на основе структуры программного обеспечения (Branch Testing и Data Flow Testing). Для каждой техники приведены проверяемые условия (test cover items), сформированные контрольные примеры, а также рассмотрены преимущества и недостатки. В заключительной части главы обоснован выбор библиотеки модульного тестирования и описана реализация тестов с использованием функционала assertion, assumption, мокирования, параметризованных тестов и матчеров.

\section{Технологии проектирования тестов на основе спецификации} \label{ch3:sec1}

Тестирование на основе спецификации (specification-based, или <<чёрный ящик>>) предполагает построение контрольных примеров исходя из описания внешнего поведения программы без учёта её внутренней реализации. Из перечня, приведённого в задании, выбраны два метода: \textbf{Boundary Value Analysis} (далее --- BVA) и \textbf{Equivalence Partitioning} (далее --- EP).

\subsection{Boundary Value Analysis (BVA)} \label{ch3:sec1:bva}

\textbf{Описание метода.} Анализ граничных значений основан на эмпирическом наблюдении: дефекты в программном обеспечении чаще всего возникают на границах диапазонов входных данных и условий ветвления. Метод заключается в выборе таких контрольных примеров, входные значения которых лежат на границе допустимой области, непосредственно перед ней и сразу после неё. Для каждой граничной точки формируется как минимум один тест.

\textbf{Применение в проекте.} Для приложения фильтрации IP-адресов выделены следующие граничные значения:
\begin{itemize}
    \item размер списка проверяемых IP-адресов: 0 элементов (пустой список), 1 элемент (минимально содержательный случай), большое количество элементов (любое число, т.к. у диапазона нет верхней границы);
    \item состояние внутреннего фильтра: пустой фильтр, фильтр с одним элементом, заполненный фильтр;
    \item корректность входной строки IP-адреса: <<типовое>> значение, пустая строка \texttt{""}, значение \texttt{null}.
\end{itemize}

\textbf{Проверяемые условия (test cover items).}
\begin{enumerate}
    \item Обработка пустого списка IP-адресов сервисом \texttt{FilterService.processAndCompare}: возвращается корректно сформированный объект \texttt{OutputData} с пустым списком результатов и заполненной статистикой.
    \item Корректная обработка одного элемента (минимальная граница содержательного входа).
    \item Корректная обработка большого числа элементов без потери данных.
    \item Поведение фильтра \texttt{HashSetFilter} и \texttt{BloomFilter} на пустом фильтре: метод \texttt{contains} возвращает \texttt{false} для любого запроса.
    \item Граничные значения для входной строки: \texttt{null} и \texttt{""} не приводят к увеличению размера фильтра и не вызывают исключений.
\end{enumerate}

\textbf{Контрольные примеры.} Сформированные тестовые случаи для метода BVA приведены в таблице~\ref{tab:bva-cases}.

\begin{table}[ht]
\centering\small
\caption{Контрольные примеры, полученные методом BVA}
\label{tab:bva-cases}
\begin{tabular}{|p{3.6cm}|p{4.0cm}|p{4.0cm}|p{2.6cm}|}
\hline
\textbf{Тестовый случай} & \textbf{Входные данные} & \textbf{Ожидаемый результат} & \textbf{Тестовый метод} \\
\hline
Пустой список IP & \texttt{[]} & \texttt{OutputData} с 0 результатами, валидной статистикой и рекомендацией & \texttt{testBoundary\_Empty\-List} \\
\hline
Один элемент (нижняя граница содержательного входа) & \texttt{["192.168.1.1"]} & 1 результат проверки, поле \texttt{ip} заполнено & \texttt{testBoundary\_Single\-Element} \\
\hline
Большое количество элементов & неограниченное кол-во различных IP & соответствующее кол-во результатов, заполненная статистика & \texttt{testBoundary\_Many\-Elements} \\
\hline
Пустой фильтр + произвольный IP & \texttt{HashSetFilter()}, запрос \texttt{"192.168.1.1"} & \texttt{contains} возвращает \texttt{false} & \texttt{testBoundary\-Values\_Empty\-And\-SingleElement} \\
\hline
Граница: добавление 1 элемента и проверка соседнего IP & \texttt{add("192.168.1.1")}, \texttt{contains("192.168.1.2")} & \texttt{true}, затем \texttt{false} & \texttt{testBoundary\-Values\_Empty\-And\-SingleElement} \\
\hline
\texttt{null} как входное значение \texttt{add} & \texttt{add(null)} & размер фильтра не меняется & \texttt{testAdd\_Null\-Value} \\
\hline
Пустая строка как входное значение \texttt{add} & \texttt{add("")} & размер фильтра не меняется & \texttt{testAdd\_Empty\-String} \\
\hline
\end{tabular}
\end{table}

\subsection{Equivalence Partitioning (EP)} \label{ch3:sec1:ep}

\textbf{Описание метода.} Разбиение на классы эквивалентности основано на предположении, что входные данные программы можно разделить на группы, в пределах которых поведение программы одинаково. Тогда достаточно протестировать по одному представителю каждого класса. Различают валидные классы (для которых ожидается корректное поведение) и невалидные (для которых ожидается отказ или специальная реакция).

\textbf{Применение в проекте.} Для входных данных метода \texttt{contains} и \texttt{add} выделены следующие классы эквивалентности:
\begin{itemize}
    \item \textit{валидные IPv4-адреса} (например, \texttt{10.0.0.1}, \texttt{172.16.0.1}, \texttt{8.8.8.8}, \texttt{192.168.1.1}) --- ожидается корректная обработка;
    \item \textit{синтаксически некорректные строки} (\texttt{invalid.ip}, \texttt{256.256.256.256}, \texttt{abc.def.ghi.jkl}) --- ожидается отсутствие исключений;
    \item \textit{пустая строка и \texttt{null}} --- ожидается, что фильтр не изменит своё состояние;
    \item \textit{дубликат уже добавленного адреса} --- ожидается, что размер не увеличивается (для \texttt{HashSetFilter}).
\end{itemize}

\textbf{Проверяемые условия (test cover items).}
\begin{enumerate}
    \item Каждый IP-адрес из класса <<валидные адреса>> добавляется и затем находится методом \texttt{contains}.
    \item Любой адрес из класса <<синтаксически некорректные>> не приводит к выбросу исключения при обработке через \texttt{processAndCompare}.
    \item Адреса из класса <<\texttt{null} и пустая строка>> не изменяют размер фильтра.
    \item Дубликаты в \texttt{HashSetFilter} не увеличивают \texttt{getSize()}.
\end{enumerate}

\textbf{Контрольные примеры.} Сформированные кейсы реализованы как параметризованные тесты, что позволяет проверить всех представителей класса одной тестовой функцией. Примеры приведены в таблице~\ref{tab:ep-cases}.

\begin{table}[ht]
\centering\small
\caption{Контрольные примеры, полученные методом Equivalence Partitioning}
\label{tab:ep-cases}
\begin{tabular}{|p{3.6cm}|p{3.8cm}|p{4.2cm}|p{2.6cm}|}
\hline
\textbf{Класс эквивалентности} & \textbf{Представители} & \textbf{Ожидаемый результат} & \textbf{Тестовый метод} \\
\hline
Валидные IPv4 & \texttt{10.0.0.1}, \texttt{172.16.0.1}, \texttt{8.8.8.8} & \texttt{contains == true} после \texttt{add} & \texttt{testEquivalence\-Partitioning\_Valid\-Ips} \\
\hline
Валидные IPv4 (для сервиса) & \texttt{192.168.1.1}, \texttt{10.0.0.1}, \texttt{172.16.0.1}, \texttt{8.8.8.8}, \texttt{1.1.1.1} & 1 результат, поле \texttt{ip} совпадает с входом & \texttt{testEquivalence\_Valid\-IpAddresses} \\
\hline
Синтаксически некорректные & \texttt{invalid.ip}, \texttt{256.256.256.256}, \texttt{abc.def.ghi.jkl} & Исключение не выбрасывается, формируется отчёт & \texttt{testEquivalence\_Invalid\-IpAddresses} \\
\hline
\texttt{null} и пустая строка & \texttt{null}, \texttt{""} & \texttt{getSize() == 0} & \texttt{testEquivalence\-Partitioning\_Invalid\-Inputs} \\
\hline
Дубликаты (HashSet) & \texttt{add("192.168.1.1")} дважды & \texttt{getSize() == 1} & \texttt{testAdd\_Duplicate\-Values} \\
\hline
\end{tabular}
\end{table}

\textbf{Оптимизация контрольных примеров.} Метод EP допускает оптимизацию по принципу <<один представитель на класс>>: вместо проверки всех возможных значений выбирается минимально достаточное множество. В результате количество тестов сокращается с потенциально неограниченного до 5 классов $\times$ 1--3 представителя $=$ 5--15 кейсов вместо тысяч возможных строковых значений.

\section{Технологии проектирования тестов на основе структуры ПО} \label{ch3:sec2}

Тестирование на основе структуры ПО (structure-based, или <<белый ящик>>) опирается на знание исходного кода и его управляющих/информационных потоков. Из списка, приведённого в задании, выбраны два метода: \textbf{Branch Testing} и \textbf{Data Flow Testing}.

\subsection{Branch Testing} \label{ch3:sec2:branch}

\textbf{Описание метода.} Тестирование ветвей предполагает построение такого набора тестов, при выполнении которого каждая ветвь управления (исход условного оператора, в т.ч. \texttt{else}-ветвь) выполняется хотя бы один раз. По сравнению с тестированием инструкций (statement testing) данный метод сильнее, поскольку гарантирует прохождение и положительной, и отрицательной ветви каждого условия.

\textbf{Применение в проекте.} В коде приложения выделены следующие критические условия:
\begin{itemize}
    \item проверка \texttt{if (ip == null)} в методах \texttt{add} и \texttt{contains} обоих фильтров;
    \item проверка \texttt{if (ip != null \&\& !ip.isEmpty())} в \texttt{HashSetFilter.add};
    \item сравнение результатов двух фильтров и установка флага \texttt{match} в \texttt{FilterService.processAndCompare};
    \item ветвление в \texttt{FilterService.generateRecommendation} в зависимости от соотношения памяти HashSet и Bloom Filter;
    \item цикл \texttt{for (i = 0; i < hashFunctionsCount; i++)} в \texttt{BloomFilter.contains} (граничная ветвь выхода из цикла).
\end{itemize}

\textbf{Проверяемые условия (test cover items).}
\begin{enumerate}
    \item Истинная и ложная ветвь \texttt{ip == null} в \texttt{contains} (обе реализации фильтров).
    \item Ветвь полного совпадения результатов двух фильтров (\texttt{match == true}).
    \item Ветвь, в которой оба фильтра вернули \texttt{false} (отрицательный сценарий, см. \texttt{testBranch\_ResultSetMatchTrue}).
    \item Ветвь успешного формирования и заполнения объекта \texttt{FilterStatistics}.
    \item Ветвь генерации непустой строки рекомендации.
\end{enumerate}

\textbf{Контрольные примеры.} Примеры реализованных тестов приведены в таблице~\ref{tab:branch-cases}.

\begin{table}[ht]
\centering\small
\caption{Контрольные примеры, полученные методом Branch Testing}
\label{tab:branch-cases}
\begin{tabular}{|p{4.2cm}|p{4.2cm}|p{3.5cm}|p{2.6cm}|}
\hline
\textbf{Покрываемая ветвь} & \textbf{Входные данные} & \textbf{Ожидаемый результат} & \textbf{Тестовый метод} \\
\hline
\texttt{ip == null} (истина) в \texttt{contains} & \texttt{contains(null)} & возвращается \texttt{false} & \texttt{testBranch\-Coverage\_Null\-Input\-Contains} \\
\hline
\texttt{ip == null} (ложь) в \texttt{contains} & \texttt{contains("1.1.1.1")} (фильтр пуст) & возвращается \texttt{false} & \texttt{testBranch\_Not\-Found\-Returns\-False} \\
\hline
\texttt{match == true} (оба фильтра вернули \texttt{false}) & \texttt{processAndCompare(["8.8.8.8"])} без загруженного blacklist & \texttt{isMatch() == true}, оба результата \texttt{false} & \texttt{testBranch\_Result\-Set\-Match\-True} \\
\hline
Ветвь заполнения статистики & \texttt{processAndCompare(["1.1.1.1"])} & все поля \texttt{FilterStatistics} $\geq 0$, вероятность FP $\in [0, 1]$ & \texttt{testBranch\_Statistics\-Fields\-Populated} \\
\hline
Ветвь генерации рекомендации & \texttt{processAndCompare(["1.1.1.1"])} & непустая строка длиной более 10 символов & \texttt{testBranch\_Recom\-mendation\-Generated} \\
\hline
\end{tabular}
\end{table}

\subsection{Data Flow Testing} \label{ch3:sec2:dataflow}

\textbf{Описание метода.} Тестирование потока данных рассматривает программу как граф, в котором отслеживаются точки определения переменных (definitions, def), их использования в вычислениях (compute use, c-use) и их использования в условиях (predicate use, p-use). Цель --- покрыть пары <<определение--использование>> (def-use pairs), чтобы каждая <<живая>> цепочка <<определение значения~$\rightarrow$~его потребление>> была проверена хотя бы одним тестом. Метод позволяет обнаруживать ошибки, связанные с некорректной инициализацией, преждевременным переопределением или потерей значения переменной.

\textbf{Применение в проекте.} Выделены следующие цепочки def-use:
\begin{itemize}
    \item \texttt{getSize()} в фильтрах: \textit{def} в конструкторе ($=0$), последовательные \textit{def} в методе \texttt{add}, \textit{c-use} в \texttt{getSize}, \textit{def} (сброс в 0) в \texttt{clear};
    \item \texttt{OutputData} в \texttt{FilterService.processAndCompare}: \textit{def} объекта, \textit{def} полей \texttt{timestamp}, \texttt{results}, \texttt{statistics}, \texttt{recommendation}, \textit{c-use} каждого из этих полей при возврате;
    \item \texttt{FilterResult.match}: \textit{def} как сравнение \texttt{hashSetResult} и \texttt{bloomFilterResult}, \textit{c-use} при выводе и при формировании отчёта.
\end{itemize}

\textbf{Проверяемые условия (test cover items).}
\begin{enumerate}
    \item Цепочка \textit{def(в конструкторе) $\to$ c-use(getSize)}: после создания фильтра размер равен нулю.
    \item Цепочка \textit{def(в add) $\to$ c-use(getSize)}: каждое добавление увеличивает счётчик.
    \item Цепочка \textit{def(в clear) $\to$ c-use(getSize)}: после \texttt{clear} счётчик снова равен нулю.
    \item Цепочка \textit{def(полей OutputData в processAndCompare) $\to$ c-use(в потребителе)}: все поля доступны и не равны \texttt{null}.
    \item Цепочка \textit{def(match в processAndCompare) $\to$ c-use(в FilterResult)}: значение \texttt{match} согласовано с \texttt{hashSetResult == bloomFilterResult}.
\end{enumerate}

\textbf{Контрольные примеры.} Сформированные кейсы приведены в таблице~\ref{tab:dataflow-cases}.

\begin{table}[ht]
\centering\small
\caption{Контрольные примеры, полученные методом Data Flow Testing}
\label{tab:dataflow-cases}
\begin{tabular}{|p{5.2cm}|p{4.2cm}|p{3.4cm}|p{3.5cm}|}
\hline
\textbf{Цепочка def--use} & \textbf{Сценарий} & \textbf{Ожидаемый результат} & \textbf{Тестовый метод} \\
\hline
\texttt{getSize}: def~$\to$~c-use & создание~$\to$~\texttt{add("A")}~$\to$~\texttt{add("B")}~$\to$~\texttt{clear} & 0~$\to$~1~$\to$~2~$\to$~0 & \texttt{testData\-Flow\_Elements\-Count\-Tracking} \\
\hline
\texttt{HashSetFilter}: def~$\to$~c-use & \texttt{new}~$\to$~\texttt{add}~$\to$~\texttt{getSize}/\texttt{contains} & \texttt{size==1}, \texttt{contains==true} & \texttt{testData\-Flow\_Definition\-Use} \\
\hline
\texttt{OutputData}: def полей $\to$ c-use & \texttt{processAndCompare(["192.168.1.1"])}, проверка всех полей & все поля не \texttt{null}, размер списка результатов $==1$ & \texttt{testData\-Flow\_Output\-Data\-Creation} \\
\hline
\texttt{FilterResult.match}: def $\to$ c-use & \texttt{processAndCompare(["10.0.0.1"])} & \texttt{isMatch() == (isHashSet\-Result() == isBloom\-FilterResult())} & \texttt{testData\-Flow\_Filter\-Result\-Creation} \\
\hline
\end{tabular}
\end{table}

\section{Сравнительный анализ методов проектирования тестов} \label{ch3:sec3}

В таблице~\ref{tab:methods-compare} приведено сравнение четырёх применённых методов проектирования тестов по числу полученных контрольных примеров, типу покрытия и трудозатратности.

\begin{table}[ht]
\centering\small
\caption{Сравнительный анализ применённых методов проектирования тестов}
\label{tab:methods-compare}
\begin{tabular}{|p{3.0cm}|p{2.0cm}|p{3.5cm}|p{3.0cm}|p{2.5cm}|}
\hline
\textbf{Метод} & \textbf{Кол-во кейсов} & \textbf{Тип покрытия} & \textbf{Преимущества} & \textbf{Недостатки} \\
\hline
Boundary Value Analysis & 7 & Граничные значения входных и внутренних диапазонов & Высокая эффективность поиска ошибок на границах & Не обнаруживает дефекты вне граничных областей \\
\hline
Equivalence Partitioning & 5 (с параметризацией покрывает 11 представителей) & Классы эквивалентности входных данных & Существенное сокращение числа тестов при сохранении покрытия & Опирается на корректное выделение классов; некорректное разбиение оставляет дефекты непокрытыми \\
\hline
Branch Testing & 5 & Все ветви управления & Гарантирует выполнение положительной и отрицательной ветви каждого условия & Не покрывает комбинации условий; требует знания исходного кода \\
\hline
Data Flow Testing & 4 & Пары def--use переменных & Выявляет ошибки инициализации и преждевременного переопределения & Высокая трудоёмкость анализа потоков данных; требует средств статического анализа \\
\hline
\end{tabular}
\end{table}

\textbf{Оценка трудозатратности.} Наименее трудозатратными в данной работе оказались методы EP и BVA: разбиение входной области и выделение границ выполнено по одному просмотру спецификации. Branch Testing потребовал ручного анализа исходного кода всех методов фильтрации и сервиса. Data Flow Testing оказался наиболее трудоёмким --- потребовалось вручную выделять def-use цепочки для критических переменных, поскольку специализированный инструмент анализа потоков данных в проекте не использовался.

\textbf{Покрытие проверяемых условий.} Совокупно четыре техники покрывают около 21 контрольного примера, и эти кейсы дополняют друг друга: BVA и EP отвечают за внешнее поведение, Branch Testing и Data Flow Testing --- за корректность внутренней реализации. Совместное применение методов из обоих классов (specification-based и structure-based) позволило достичь высоких метрик покрытия (см. главу~\ref{ch4}).

\section{Выбор библиотеки модульного тестирования} \label{ch3:sec4}

В качестве основной библиотеки модульного тестирования выбран \textbf{JUnit~5} (Jupiter). Дополнительно использованы \textbf{Hamcrest} (для декларативных матчеров) и \textbf{Mockito} (для мокирования зависимостей). Сравнение JUnit~5 с альтернативами и обоснование выбора приведено в таблице~\ref{tab:lib-compare}.

\begin{table}[ht]
\centering\small
\caption{Сравнение библиотек модульного тестирования для Java}
\label{tab:lib-compare}
\begin{tabular}{|p{2.4cm}|p{4.2cm}|p{4.2cm}|p{3.0cm}|}
\hline
\textbf{Библиотека} & \textbf{Преимущества} & \textbf{Недостатки} & \textbf{Решение} \\
\hline
JUnit~5 & Современная архитектура (Jupiter+Vintage), модульность, поддержка параметризованных тестов, динамических тестов, расширений, штатные \texttt{Assumptions} & Более крупный набор зависимостей, чем у JUnit~4 & \textbf{Выбран} \\
\hline
JUnit~4 & Широкое распространение, простой API & Устаревшая архитектура, отсутствие нативных параметризованных тестов и assumptions в современном виде & Не выбран \\
\hline
TestNG & Гибкие группы тестов, наследование конфигурации & Меньшая интеграция в инструменты сборки и Pitest на этапе работы & Не выбран \\
\hline
\end{tabular}
\end{table}

\textbf{Основные используемые команды и аннотации.}
\begin{itemize}
    \item \texttt{@Test} --- разметка тестового метода;
    \item \texttt{@BeforeEach} --- инициализация фильтра и/или сервиса перед каждым тестом;
    \item \texttt{@ParameterizedTest} с \texttt{@ValueSource} и \texttt{@CsvSource} --- параметризованные тесты для классов эквивалентности;
    \item \texttt{Assertions.assertEquals}, \texttt{assertTrue}, \texttt{assertFalse}, \texttt{assertNotNull}, \texttt{assertAll}, \texttt{assertDoesNotThrow}, \texttt{assertThrows} --- семейство ассертов;
    \item \texttt{Assumptions.assumeTrue}, \texttt{assumeFalse}, \texttt{assumingThat} --- предположения для условного выполнения проверок;
    \item \texttt{org.hamcrest.MatcherAssert.assertThat} с матчерами \texttt{equalTo}, \texttt{greaterThan}, \texttt{lessThan}, \texttt{notNullValue}, \texttt{not(emptyString())} --- декларативная проверка;
    \item \texttt{Mockito.mock}, \texttt{when(...).thenReturn(...)}, \texttt{verify}, \texttt{spy}, \texttt{doThrow} --- работа с заглушками и шпионами;
    \item \texttt{mvn test} --- запуск всех тестов через Maven Surefire.
\end{itemize}

\section{Реализация тестов и используемый функционал JUnit\,5} \label{ch3:sec5}

Тесты разработаны по принципам TDD: для каждого нового публичного метода сначала писался падающий тест, затем минимальная реализация, затем рефакторинг. Итоговый набор включает $\sim$120 тестовых методов, распределённых по пяти классам тестов: \texttt{HashSetFilterTest}, \texttt{BloomFilterTest}, \texttt{FilterServiceTest}, \texttt{FilterServiceBoundaryTest}, \texttt{FilterServiceEdgeCaseTest}.

\subsection{Использование функционала assertion (минимум 5 методов)}
В проекте применены следующие методы класса \texttt{Assertions} (демонстрационный набор приведён в \texttt{FilterServiceTest.testAssertionsVariety}):
\begin{enumerate}
    \item \texttt{assertEquals(expected, actual)} --- проверка равенства значений;
    \item \texttt{assertTrue(condition)} --- проверка истинности условия;
    \item \texttt{assertFalse(condition)} --- проверка ложности условия;
    \item \texttt{assertNotNull(object)} --- проверка ненулевой ссылки;
    \item \texttt{assertAll("group", ...)} --- группировка нескольких ассертов с накоплением ошибок;
    \item \texttt{assertDoesNotThrow(executable)} --- проверка отсутствия исключений;
    \item \texttt{assertThrows(IOException.class, executable)} --- проверка ожидаемого исключения (используется в тестах с моком \texttt{JsonHandler}).
\end{enumerate}
Итого использовано не менее семи различных методов assertion, что превышает требование задания (минимум~5).

\subsection{Использование функционала assumption (минимум 2 метода)}
Применены три метода класса \texttt{Assumptions}:
\begin{enumerate}
    \item \texttt{assumeTrue(condition, "сообщение")} --- если условие не выполнено, тест помечается как пропущенный (используется в \texttt{testAssumptionsUsage} для проверки, что сервис создан);
    \item \texttt{assumeFalse(condition, "сообщение")} --- симметричный вариант (используется для проверки, что список результатов не пуст);
    \item \texttt{assumingThat(condition, executable)} --- условное выполнение блока проверок (используется в \texttt{testAssumptions\_Environment} класса \texttt{BloomFilterTest}).
\end{enumerate}

\subsection{Мокирование (минимум 3 вида)}
В классе \texttt{FilterServiceTest} применены четыре сценария работы с библиотекой Mockito:
\begin{enumerate}
    \item \textbf{Стандартный mock с настройкой возвращаемого значения} --- \texttt{mock(IpFilter.class)} с \texttt{lenient().when(...).thenReturn(...)} (тест \texttt{testMockBehavior\_VerifyWithLenient});
    \item \textbf{Verify взаимодействий} --- \texttt{verify(mock, times(n)).method(...)} и \texttt{verify(mock, never()).method(...)} (тест \texttt{testMockVerify\_Interactions});
    \item \textbf{Spy (частичный мок)} --- \texttt{spy(new ArrayList<>(...))} с проверкой как вызовов мока, так и реальных значений (тест \texttt{testMockSpy\_Usage});
    \item \textbf{Mock, выбрасывающий исключение} --- \texttt{doThrow(new IOException(...)).when(mock).method(...)} в комбинации с \texttt{assertThrows} (тест \texttt{testMockException\_Throwing}).
\end{enumerate}
Указанные сценарии иллюстрируют различные аспекты работы с зависимостями: подмена результата, проверка взаимодействия, частичная подмена и эмуляция отказов.

\subsection{Параметризованные тесты}
Для реализации классов эквивалентности и серий граничных случаев использованы параметризованные тесты JUnit~5:
\begin{itemize}
    \item \texttt{@ParameterizedTest} + \texttt{@ValueSource(strings = \{...\})} --- проверка набора валидных IPv4 (\texttt{testEquivalence\-Partitioning\_Valid\-Ips}, \texttt{testEquivalence\_Valid\-IpAddresses});
    \item \texttt{@ParameterizedTest} + \texttt{@ValueSource(ints = \{...\})} --- проверка отклика сервиса на разный объём входных данных (\texttt{testParameterized\_DifferentIp\-Counts});
    \item \texttt{@ParameterizedTest} + \texttt{@CsvSource(...)} --- комбинированные пары входных данных и ожидаемого результата (\texttt{testEquivalence\-Classes}, \texttt{testEquivalence\_Invalid\-IpAddresses}).
\end{itemize}

\subsection{Матчеры (минимум 2 вида)}
Для декларативных проверок использована библиотека Hamcrest. В тесте \texttt{testHamcrestMatchers} класса \texttt{FilterServiceTest} применены следующие матчеры:
\begin{enumerate}
    \item \texttt{greaterThanOrEqualTo} --- сравнение числовых значений памяти;
    \item \texttt{lessThanOrEqualTo} --- проверка верхней границы вероятности;
    \item \texttt{equalTo} --- сравнение значений;
    \item \texttt{notNullValue} --- проверка ненулевой ссылки;
    \item \texttt{not(emptyString())} --- проверка непустой строки рекомендации.
\end{enumerate}
Дополнительно в \texttt{BloomFilterTest} применены \texttt{lessThan} и \texttt{greaterThanOrEqualTo}.

\subsection{Сводка по требованиям задания}
В таблице~\ref{tab:reqs-checklist} приведена сверка реализованного функционала с количественными требованиями задания.

\begin{table}[ht]
\centering\small
\caption{Соответствие реализации требованиям задания к JUnit-функционалу}
\label{tab:reqs-checklist}
\begin{tabular}{|p{4.5cm}|p{2.5cm}|p{2.5cm}|p{4.0cm}|}
\hline
\textbf{Требование} & \textbf{Минимум} & \textbf{Реализовано} & \textbf{Где} \\
\hline
Методы assertion & 5 & 7 & \texttt{testAssertionsVariety} \\
\hline
Методы assumption & 2 & 3 & \texttt{testAssumptionsUsage}, \texttt{testAssumptions\_Environment} \\
\hline
Виды мокирования & 3 & 4 & \texttt{testMockBehavior\_*}, \texttt{testMockVerify\_*}, \texttt{testMockSpy\_*}, \texttt{testMockException\_*} \\
\hline
Параметризованные тесты & 1+ & 4 параметризованных метода & \texttt{testParameterized\_*}, \texttt{testEquivalence*} \\
\hline
Матчеры & 2 & 5 & \texttt{testHamcrestMatchers} \\
\hline
\end{tabular}
\end{table}

\subsection{Тестирование основного алгоритма по принципу сравнения с альтернативой}
В соответствии с требованием задания, основной наукоёмкий алгоритм (вероятностный Bloom Filter) тестируется в том числе путём сравнения его результатов с альтернативной реализацией (детерминированный HashSet). Метод \texttt{FilterService.processAndCompare} последовательно вызывает \texttt{contains} обеих реализаций, формирует поле \texttt{match} и фиксирует расхождения. Тесты \texttt{testBranch\_ResultSetMatchTrue} и \texttt{testDataFlow\_FilterResultCreation} верифицируют согласованность двух алгоритмов, а допустимые расхождения интерпретируются как False Positive со стороны Bloom Filter (это поведение явно проверяется в \texttt{testNoFalseNegatives}: добавленный элемент обязан находиться обоими алгоритмами).

\newpage
